\sscp{\tsc{Teun-Nila-Serua}}{04-06-04}
\largerpage
The Teun, Nila and Serua languages are each spoken on three islands
of the same names to the northeast of Damer Island.\footnote{
		The linguistic situation in Teun,
		Nila, and Serua was actually more complex.
		Teun was also spoken in the village of Bumei on Nila Island.
		There were also two Wetan speaking villages,
		Isu and Layeni, in southeastern Teun Island.
		\citet{Chlenova2006} provides data for this Wetan variety.
		\citet[52, 78]{vanEngelenhoven2003}, who appears to have worked
		only with Serua consultants,
		reports that Serua was the lingua franca used between inhabitants
		of each island and on sea voyages,
		and that Serua was spoken in Sifluru
		and Wotay in south Nila.
		Edwards’ Nila consultants reported that Nila was spoken in all
		villages of Nila Island (except Bumei; Teun speaking).
		Edwards’ data from Wotay shows that this is indeed the Nila language,
		with only slight differences from that in \citet{Taber1993}.}
From 1976--1983 the speakers were relocated to southern Seram,
but there are also reports as early as \citet[404]{Taber1993}
that some speakers have moved back to their home islands.
\citet{Edwards2023} has 505-item word lists for each language.
Apart from this, \citet{Taber1993} provides 193-item word lists
for all three languages.\footnote{
		There is a Holle list collected
		in 1937 and published in \citet[161--174]{Stokhof1981b} labelled “Niala”
		which is “from the village of Boemei [Bumei]”.
		Comparison of this list with other Teun data shows that this
		list represents the Teun language,
		which was indeed spoken in Bumei village on Nila Island (see the preceding footnote).
		This list was reportedly collected in “Latoerake [Laturake],
		Mountains in the Riring region, West-Ceram”.
		The origins of this list represent an unresolved puzzle.
		\citet[197]{Collins1982d} states regarding this list (based on
		the limited metadata which accompanies it): “The informant,
		Matheus Soeplatoe, was a native speaker of Alune; he was the
		head of a village near Riring,
		the military garrison and prison of West Seram in the colonial period.
		[{\ldots}] How Bpk Matheus learned Nila [sic.] remains a puzzle.
		He may have served as a soldier or catechist there or he may
		have learned the language through a prisoner in the forced labour
		camp in Riring.”}
A 500-item word list for Serua is provided by \citet{Chlenova2004}.
\citet{vanEngelenhoven2003} provides some grammatical notes on Serua.
Data is from \citet{Edwards2023}, unless otherwise noted.
This data is representative of the following villages: Yafila
(Teun), Wotay (Nila), and Lesluru (Serua).

\citet[404f]{Taber1993} places Teun, Nila,
and Serua in a single group based on lexical similarity with
Nila and Serua closer to one another than either is to Teun.
Teun averages 69{\%} similarity to Nila
and 66{\%} similarity to Serua,
while Nila and Serua average 76{\%} similarity.
Phonological support for placing these languages together comes
from shared *z > \it{s},
as shown in \trf{04-34}.

\begin{table}\tcp{\tsc{Teun-Nila-Serua} *z}{04-34}
	\begin{threeparttable}\st{6pt}
	\begin{tabular}{lllll}\lsptoprule
local gl.	&	rain	&	point at	&	ladder	&	path\\
PMP	&	*qu\tbr{z}an	&	*tu\tbr{z}uq	&	*haRə\tbr{z}an	&	*\tbr{z}alan\\ \midrule
Teun	&	{\hr}u\tbr{s}na	&	{\hr}tu\tbr{s}	&	\cg{(kotna)}	&	{\hr}talla\su{†}\\
Nila	&	{\hr}u\tbr{s}na	&	{\hr}tu\tbr{s}na	&	{\hr}re\tbr{s}na	&	{\hr}\tbr{s}alna\\
Serua	&	{\hr}u\tbr{s}na	&	{\hr}n-tu\tbr{s}an	&	{\hr}re\tbr{s}na	&	{\hr}\tbr{s}alna\\
		\lspbottomrule
	\end{tabular}
		\begin{tablenotes}\tnsize
			\item[†] {Teun \it{talla} is almost certainly borrowed from Wetan \it{talna}.
								Taber has Teun \it{kala} with unexplained initial \it{k}
								the closest match for which is Kisar \it{kalla} ‘path’.}
		\end{tablenotes}
	\end{threeparttable}
\end{table}

Additional support for \tsc{Teun-Nila-Serua} (TNS) comes from
shared *aw > \it{a} (\trf{04-35}) and,
perhaps, *R/*j > \it{r} (see below).
However, both these changes are also found in \tsc{Kisar-Luang},
and East Damer and may instead provide evidence for a higher
level node containing these groups.
Additional support for TNS,
which separates it from other languages of \tsc{Southwest Maluku},
comes from shared *ŋ > \it{k} in *laŋit > \it{lakti} ‘sky’ which
is unique to TNS.

Within TNS, there is evidence for a \tsc{Teun-Nila} node on the basis
of shared phonological changes.
This conflicts with the evidence from lexical similarity which
places Nila and Serua together.
The first two changes shared by Teun
and Nila are *ay > \it{a},
and *b > **f.
The variety of Nila in \citet{Taber1993},
from Kokroman village (represented as “Nila (K)”),
has subsequent initial **f > \it{f,
h} and intervocalic **f > \it{h.} (The split of initial *b- >
**f- > \it{h-, f-} in Kokroman Nila is roughly even in \citet{Taber1993} with
about half a dozen instances each.) Serua,
on the other hand,
has *ay > \it{i}
and usually has *b > \it{w} with subsequent **w > {\0} /V{\gap}\it{u}.
Examples of *ay are given in \trf{04-35}
and examples of *b are given in \trf{04-36}.

\begin{table}\tcp{\tsc{Teun-Nila-Serua} *-ay and *-aw}{04-35}
\begin{adjustbox}{width=\textwidth}
	\st{2pt}\begin{tabular}{lllll|llll}\lsptoprule
local gl.	&	man, male	&	die	&	sand	&	liver	&	mouse	&	day	&	steal	&	rafter\\
PMP	&	*maRuqan\tbr{ay}	&	*mat\tbr{ay}	&	*qən\tbr{ay}	&	*qat\tbr{ay}	&	*balab\tbr{aw}	&	*qaləj\tbr{aw}	&	*nak\tbr{aw}	&	*kas\tbr{aw}\\ \midrule
Teun	&	{\hr}m-man\tbr{a}	&	{\hr}n-maʔ\tbr{a}	&	{\hr}wen\tbr{a}	&	{\hr}waʔ\tbr{a}	&	{\hr}flaf\tbr{a}	&	{\hr}low\tbr{a}	&	{\hr}na-mna\tbr{a}	&	{\hr}as\tbr{a}\\
%\multirow{2}{*}{$\left.\hspace{-2mm}\begin{array}{l}{\textrm{Nila,}}\\{\textrm{Wotay}}\\\end{array}\hspace{-1.3mm}\right\}$} &
%\multirow{2}{*}{{\hr}mon\tbr{a}}	&	\multirow{2}{*}{{\hr}n-mat\tbr{a}}	&	\multirow{2}{*}{{\hr}wen\tbr{a}}	&	\multirow{2}{*}{{\hr}y-at\tbr{a}}	&
%\multirow{2}{*}{{\hr}flaf\tbr{a}}	&	\multirow{2}{*}{{\hr}ler\tbr{a}}	&	\multirow{2}{*}{{\hr}mna\tbr{a}}	&	\multirow{2}{*}{{\hr}as\tbr{a}} \\
%&&&&&&&&\\
%\multirow{2}{*}{$\left.\hspace{-2mm}\begin{array}{l}{\textrm{Nila,}}\\{\textrm{Kokroman}}\\\end{array}\hspace{-1.3mm}\right\}$} &
%\multirow{2}{*}{{\hr}monr\tbr{a}}	&	\multirow{2}{*}{{\hr}n-mat\tbr{a}}	&	\multirow{2}{*}{{\hr}wen\tbr{a}}	&	\multirow{2}{*}{{\hr}y-at\tbr{a}}	&
%\multirow{2}{*}{{\hr}flah\tbr{a}}	&		&		&	\\
%&&&&&&&&\\
Nila (W)	&	{\hr}mon\tbr{a}	&	{\hr}n-mat\tbr{a}	&	{\hr}wen\tbr{a}	&	{\hr}y-at\tbr{a}	&	{\hr}flaf\tbr{a}	&	{\hr}ler\tbr{a}	&	{\hr}mna\tbr{a}	&	{\hr}as\tbr{a}\\
%Wotay&&&&&&&&\\
Nila (K)	&	{\hr}monr\tbr{a}	&	{\hr}n-mat\tbr{a}	&	{\hr}wen\tbr{a}	&	{\hr}y-at\tbr{a}	&	{\hr}flah\tbr{a}	&		&		&	\\
%Kokroman &&&&&&&&\\
Serua	&	{\hr}man\tbr{i}	&	{\hr}n-mat\tbr{i}	&	{\hr}en\tbr{i}	&	{\hr}n-et\tbr{i}	&	{\hr}law\tbr{a}	&	{\hr}ler\tbr{a}	&	{\hr}ne-mna\tbr{a}	&	\\
		\lspbottomrule
	\end{tabular}
\end{adjustbox}
\end{table}

\begin{table}\tcp{\tsc{Teun-Nila-Serua} *b}{04-36}
\begin{adjustbox}{width=\textwidth}
	\begin{threeparttable}\st{2pt}
\begin{tabular}{llllllllll}\lsptoprule
local gl.	&	fruit	&	kill	&	moon	&	stone	&	pig	&	mouse	&	\it{Derris}	&	soil	&	sugarcane\\
PMP	&	*\tbr{b}uaq	&	*\tbr{b}unuq	&	*\tbr{b}ulan	&	*\tbr{b}atu	&	*\tbr{b}a\tbr{b}uy	&	*\tbr{b}ala\tbr{b}aw	&	*tu\tbr{b}ah	&	*qa\tbr{b}u	&	*tə\tbr{b}uh\\ \midrule
Teun	&	{\hr}\tbr{f}ua	&	{\hr}\tbr{f}unu	&	{\hr}\tbr{f}ulla	&	{\hr}\tbr{f}aʔu	&	{\hr}\tbr{f}a\tbr{f}i	&	{\hr}\tbr{f}la\tbr{f}a	&	{\hr}tu\tbr{f}a	&	\cg{(rau)}\su{†}	&	{\hr}te\tbr{f}u\\
Nila (W)	&	{\hr}\tbr{f}ua	&	{\hr}\tbr{f}unu	&	{\hr}\tbr{f}ulna	&	{\hr}\tbr{f}atu	&	{\hr}\tbr{f}a\tbr{f}i	&	{\hr}\tbr{f}la\tbr{f}a	&	{\hr}tu\tbr{f}a	&	{\hr}wa\tbr{f}u	&	{\hr}to\tbr{f}u\\
Nila (K)	&	{\hr}m-\tbr{f}ua	&	{\hr}\tbr{f}ulnu	&	{\hr}\tbr{h}ulna	&	{\hr}\tbr{h}atu	&		&	{\hr}\tbr{f}la\tbr{h}a	&		&	{\hr}wa\tbr{h}u	&	\\
Serua	&	{\hr}\tbr{w}ua	&	{\hr}\tbr{w}unu	&	{\hr}\tbr{w}ulna	&	{\hr}\tbr{w}atu	&	{\hr}\tbr{w}a\tbr{w}i	&	{\hr}la\tbr{w}a	&	{\hr}tu\tbr{w}a	&	{\hr}au	&	{\hr}tou\\
		\lspbottomrule	\end{tabular}
		\begin{tablenotes}\tnsize
			\item[†] {\citet{Edwards2023} also has Nila [ʔafu] ‘ash’,
								which may be a loan given initial *q > {\0} (expect subsequent
								*{\0} > \it{w}, see below).
								\citet{Edwards2023} also has Teun \it{lafu} ‘ash’ while \citet[35]{Taber1993}
								has Teun \it{yafu} each apparently from earlier **liabu,
								as attested by East Damer \it{liawu}
								and North Babar \it{niawi}.}
		\end{tablenotes}
	\end{threeparttable}
\end{adjustbox}
\end{table}

A third change shared by Teun
and Nila is insertion of the glide \it{w} before vowel initial words.
This change is more complex than the other two,
and requires detailed discussion,
as well as consideration of the reflexes of PMP *w.
Teun and Nila have retention of *w = \it{w},
as well as insertion of \it{w} before most vowel initial words.
Serua, on the other hand has no such insertion but has also lost word-initial *w.
Examples of each change are given in \trf{04-37a}
and \trf{04-37b}.

%%%%%%%%%%%%%%%%%%%%%%%%%%%%%%%%%%%%%%%%%%%%%%%%%%%%%%%%%%%%%%%%%%%%%%%%%%%%%%%%%%%%%%%%%%%%%%%%%
%%%%	addition of a table increments the table labelling by +1 in the PDF compared with code %%%%
%%%%%%%%%%%%%%%%%%%%%%%%%%%%%%%%%%%%%%%%%%%%%%%%%%%%%%%%%%%%%%%%%%%%%%%%%%%%%%%%%%%%%%%%%%%%%%%%%

\begin{table}\tcp{\tsc{Teun-Nila-Serua} *w /{\#\gap}}{04-37a}
	\begin{threeparttable}\st{6pt}
	\begin{tabular}{lllll}\lsptoprule
local gl.	&	water	&	eight	&	root	&	ySib	\\
PMP	&	*\tbr{w}ahiyəR	&	*\tbr{w}alu	&	*\tbr{w}akaR\su{†}	&	**\tbr{w}aji\su{‡}	\\ \midrule
Teun	&	{\hr}\tbr{w}eru	&	{\hr}\tbr{w}alu	&	{\hr}\tbr{w}owra	&	{\hr\hr}\tbr{w}aʔi		\\
Nila	&	{\hr}\tbr{w}eru	&	{\hr}\tbr{w}alu	&	{\hr}\tbr{w}\<y\>ara	&	{\hr\hr}\tbr{w}arni		\\
Serua	&	\hp{*w}eru	&	\hp{*w}alu	&	{\hr}n-era	&	\hp{**w}arni		\\
		\lspbottomrule
	\end{tabular}
		\begin{tablenotes}\tnsize
			\item[†]	{Nila \it{w\<y\>ara} ‘root’ is from **(n)i-wara with metathesis
								of the vowel of the prefix
								and initial consonant of the stem.
								Serua \it{n-era} ‘root’ is from intermediate **ni-ara with \tsc{3sg}
								possessive prefix and the vowel sequence **ia > \it{e}.}
			\item[‡]	{**waji is intermediate from PMP *huaji ‘younger sibling’.}
		\end{tablenotes}
	\end{threeparttable}
\end{table}

\begin{table}\tcp{\tsc{Teun-Nila-Serua} *V /{\#\gap}}{04-37b}
\begin{adjustbox}{width=\textwidth}
	\begin{threeparttable}\st{2pt}
	\begin{tabular}{llllllllll}\lsptoprule
local gl.	&	dog	&	father	&	\tsc{1sg}	&	wind	&	fire	&	calcium	&	\tsc{2sg}	&	sand	&	vein	\\
PMP	&	*asu	&	*ama	&	*aku	&	*aŋin	&	*hapuy	&	*qapuR	&	*kahu	&	*qənay	&	*uRat	\\ \midrule
Teun	&	{\hr}\tbr{w}asu	&	{\hr}\tbr{w}ama	&	{\hr}\tbr{w}aa\su{Δ}	&	{\hr}\tbr{w}anni	&	{\hr}\tbr{w}ai	&	{\hr}\tbr{w}ari	&	{\hr}\tbr{w}oo	&	{\hr}\tbr{w}ena	&	{\hr}\tbr{w}urʔa	\\
Nila	&	{\hr}\tbr{w}asu	&	{\hr}\tbr{w}am-ku	&	{\hr}\tbr{w}aʔu	&	{\hr}\tbr{w}anni	&	{\hr}\tbr{w}ai	&	{\hr}\tbr{w}aru	&	{\hr}\tbr{w}o	&	{\hr}\tbr{w}ena	&	{\hr}\tbr{w}urta	\\
Serua	&	\hp{*w}asu	&	\hp{*w}am-na	&	\hp{*w}aʔu	&	\cg{(motra)}	&	\hp{*w}ai	&	\hp{*w}aru	&	\hp{*w}oa	&	\hp{*w}enni	&	\hp{*w}urta	\\
		\lspbottomrule
	\end{tabular}
		\begin{tablenotes}\tnsize
			\item[†] {\citet[413]{Taber1993} has Teun \it{wau} ‘\tsc{1sg}’.}
		\end{tablenotes}
	\end{threeparttable}
\end{adjustbox}
\end{table}

\newpage
There are logically three explanations for this data:

\begin{itemize}
	\item[i.] {\it{w} was inserted in Teun and Nila.
						Serua has independent *w > {\0}.}
	\item[ii.] {*w was lost at the level of Proto-TNS.
							Teun and Nila subsequently inserted \it{w} before all vowel initial
							words, including those with earlier *w > {\0}.}
	\item[iii.] {Insertion of \it{w} occurred in Proto-TNS.
							Serua subsequently had *w > {\0},
							which affected both PMP *w
							and newly acquired Proto-TNS **w.}
\end{itemize}

Careful examination of the data indicates that hypothesis
(i.) or (ii.) is correct
and that insertion of \it{w} only affected Teun and Nila.
It can thus be identified as an additional sound change shared
by these two languages.
Furthermore, the data indicates that this change either spread
by diffusion from Teun into Nila or was a change in progress
at the level of Proto-Teun-Nila which was then carried out to
a greater extent in Teun.
The evidence for this scenario is twofold.

Firstly, Teun also has insertion of \it{w} intervocalically before the vowel \it{a}.
Such \it{w} insertion also occurs after loss of medial consonants,
including cases of *R > **r > {\0}.
Examples are given in \trf{04-38}.

\begin{table}\tcp{Teun {\0} > \it{w} /V{\gap}\it{a}}{04-38}
%\begin{adjustbox}{width=\textwidth}
	\begin{threeparttable}\st{3pt}
	\begin{tabular}{lllllllll}\lsptoprule
local gl.	&	faeces	&	ginger	&	fathom	&	blood	&	banana	&	root	&	sail	&	coconut\\
PMP	&	*taqi	&	*laqia	&	*dəpa	&	*daRaq	&	**biak	&	*wakaR	&	*layaR	&	*niuR\\ \midrule
Teun	&	{\hr}te\tbr{w}a	&	{\hr}li\tbr{w}a	&	{\hr}re\tbr{w}a	&	{\hr}ra\tbr{w}a	&	{\hr}fi\tbr{w}a	&	{\hr}wo\tbr{w}ra\su{‡}	&	{\hr}la\tbr{w}a	&	{\hr}no\tbr{w}a\\
Nila	&	{\hr}tea	&	{\hr}lia	&	{\hr}rea	&	{\hr}rara	&	{\hr}fia	&	{\hr}w\<y\>ara	&	{\hr}laru	&	{\hr}nuru\\
Serua	&	{\hr}tea	&		&	{\hr}rea	&	{\hr}rara	&	{\hr}bia\su{†}	&	{\hr}n-era	&	{\hr}laru	&	{\hr}nuru\\
		\lspbottomrule
	\end{tabular}
		\begin{tablenotes}\tnsize
			\item[†]	{Leslurua Serua has \it{bia} ‘banana’ \citep{Edwards2023} Waru/Terana
								Serua has \it{wia} \citep[412]{Taber1993}.
								\citet{Chlenova2004} has Serua \it{hia} ‘banana’.}
			\item[‡]	{Teun \it{wowra} ‘root’ is probably a metathesised form of **wowar.}
		\end{tablenotes}
	\end{threeparttable}
%\end{adjustbox}
\end{table}

Secondly, Teun has instances of word-initial glide insertion
in environments where Nila does not.
The variety of Teun (Mesa or Bumei village) in \citet{Taber1993}
usually has insertion of \it{y} before words with initial *i.\footnote{
		The equivalent Teun words in \citet{Edwards2023} are \it{ina} ‘mother’,
		\it{hiru} ‘how much’, and \it{hitu} ‘seven’.
		Initial \it{h} appears to be an insertion after *p > {\0},
		as there are other words in Edwards’ data which have optional
		insertion of \it{h} /{\#}i{\gap}.
		Examples include *ia > \it{ii {\tl} hii} \tsc{‘3sg’},
		*kita > \it{ita {\tl} hita} \tsc{‘1pl.incl’},
		and *si-ida > \it{hira {\tl} ira {\tl} ir} \tsc{‘3pl’}.}
Examples are *ina > Teun \it{yina},
Nila, Serua \it{ina} ‘mother’, *pija > **hija > Teun \it{yiru},
Nila, Serua \it{ira} ‘how much’,
*pitu > \it{yitu}, Nila,
Serua \it{itu} ‘seven’.\footnote{
		One word with initial \it{i}
		in Teun which does not have glide insertion
		is *hikan > \it{ina} ‘fish’.}
There are also some words with initial vowels other
than \it{i} which have insertion of \it{y}: **ai > Teun \it{yai},
Nila \it{ai} ‘boat’ (Serua \it{wroa} ‘boat’)
and **amat > Teun \it{yamʔa},
Nila, Serua \it{amta} ‘shark’, as well as one word with initial
\it{i} which has insertion of \it{w} in Teun,
but not Nila: *quliŋ > Teun \it{willi},
Nila \it{ilni}, Serua \it{ini} ‘rudder’.

This data can be explained if glide insertion was a change which
only affected Teun and Nila and that it was more thorough
and continued in a greater range of environments in Teun than Nila.
There are thus three changes shared by Teun
and Nila: *ay > \it{i},
*b > **f, *{\0} > \it{w} /{\#}{\gap}V.

One change that sets Serua apart from Teun
and Nila is *ŋ > \it{n},
{\0} while Teun and Nila have universal
*ŋ > \it{n}, though there
are only two clear examples of *ŋ > {\0}
in the available Serua data.
Examples of *ŋ in Teun,
Nila, and Serua are given in \trf{04-39}.

\begin{table}\tcp{\tsc{Teun-Nila-Serua} *ŋ}{04-39}
\begin{adjustbox}{width=\textwidth}
	\begin{threeparttable}\st{2pt}
	\begin{tabular}{lllllllll}\lsptoprule
local gl.	&	name	&	swim	&	ear	&	cold	&	rudder	&	near	&	branch	&	soul\\
PMP	&	*\tbr{ŋ}ajan	&	*na\tbr{ŋ}uy	&	*təli\tbr{ŋ}a	&	*ma-di\tbr{\tbr{ŋ}}di\tbr{ŋ}	&	*quli\tbr{\tbr{ŋ}}	&	*da\tbr{\tbr{ŋ}}i	&	*sa\tbr{ŋ}a	&	*suma\tbr{\tbr{ŋ}}əd\\ \midrule
Teun	&	{\hr}n-\tbr{n}ana	&	{\hr}na\tbr{n}i	&	{\hr}li\tbr{n}a	&	\cg{(motra)}	&	{\hr}wil\tbr{l}i\su{†}	&	\cg{(nakatuli)}	&	{\hr}n-sa\tbr{n}a	&	{\hr}m-ma\tbr{n}a\\
Nila	&	{\hr}\tbr{n}arna	&	{\hr}na\tbr{n}i	&	{\hr}kli\tbr{n}a	&	{\hr}meri\tbr{n}a	&	{\hr}il\tbr{n}i	&	{\hr}kra\tbr{n}i	&	{\hr}sa\tbr{n}na	&	\cg{(niʔa)}\\
Serua	&	{\hr}arna	&	{\hr}nai	&	{\hr}kli\tbr{n}a	&	{\hr}mori\tbr{n}a	&	{\hr}i\tbr{n}i	&	{\hr}nikre\tbr{n}i	&	\cg{(ain tenu)}	&	\\
		\lspbottomrule
	\end{tabular}
		\begin{tablenotes}\tnsize
			\item[†] {Teun \it{willi} ‘rudder’ has assimilation of \it{n → l} after metathesis.
								Compare *bulan > \it{fulla} ‘moon’ (Nila \it{fulna,} Serua \it{wulna})
								and *zalan > \it{talla} ‘path’ (Nila,
								Serua \it{salna}).}
		\end{tablenotes}
	\end{threeparttable}
\end{adjustbox}
\end{table}

\begin{table}\tcp{\tsc{Teun-Nila-Serua} *d > \it{r}}{04-40}
\begin{adjustbox}{width=\textwidth}
	\st{2pt}\begin{tabular}{lllllllll}\lsptoprule
local gl.	&	inside	&	two	&	dugong	&	post, pole	&	cold	&	live	&	\tsc{3pl}	&	maize\\
PMP	&	*\tbr{d}aləm	&	*\tbr{d}uha	&	*\tbr{d}uyuŋ	&	*ha\tbr{d}iRi	&	*ma-\tbr{d}iŋdiŋ	&	*ma-qu\tbr{d}ip	&	*si-i\tbr{d}a	&	*bata\tbr{d}\\ \midrule
Teun	&	{\hr}\tbr{r}alma	&	{\hr}\tbr{r}uu	&	{\hr}\tbr{r}uni	&	{\hr}\tbr{r}ou	&	\cg{(motra)}	&	{\hr}mo\tbr{r}i	&	{\hr}i\tbr{r}a	&	{\hr}fet\tbr{r}a\\
Nila	&	{\hr}\tbr{r}alna	&	{\hr}\tbr{r}ua	&	{\hr}\tbr{r}uni	&	{\hr}\tbr{r}iri	&	{\hr}me\tbr{r}ina	&	{\hr}mo\tbr{r}i	&	{\hr}i\tbr{r}	&	{\hr}fet\tbr{r}a\\
Serua	&	{\hr}\tbr{r}alna	&	{\hr}\tbr{r}ua	&	{\hr}\tbr{r}uni	&	{\hr}\tbr{r}iri	&	{\hr}mo\tbr{r}ina	&	{\hr}mo\tbr{r}i	&	{\hr}i\tbr{r}a	&	{\hr}wet\tbr{r}a\\
		\lspbottomrule
	\end{tabular}
\end{adjustbox}
\end{table}

The reflexes of *d,
*R, and *j provide another perspective on the internal relationships
of Teun, Nila, and Serua.
In all three languages *d > \it{r}.
Examples are given in \trf{04-40}.
See also *dəpa ‘fathom’
and *daRaq ‘blood’ in \trf{04-38}.

PMP *j/*R > \it{r} in Nila and Serua.
Thus, these languages merge *d/*j/*R > \it{r}.
In Teun *j and *R each undergo a split to {\0}
and \it{r}, without clear conditioning environments.
Examples of *j are given in \trf{04-41}.
Another possible example of *j is *bukij > Teun,
Nila \it{fura}, Serua \it{wura} ‘mountain’.

\begin{table}\tcp{\tsc{Teun-Nila-Serua} *j}{04-41}
\begin{adjustbox}{width=\textwidth}
	\begin{threeparttable}\st{2pt}
	\begin{tabular}{lllllllll}\lsptoprule
	&		&		&		&	umbilical	&	&	&	&	\\
local gl.	&	name	&	sun	&	day	&	cord	&	ySib	&	fly	&	gall	&	how m.?\\
PMP	&	*ŋa\tbr{j}an	&	*qalə\tbr{j}aw	&	*qalə\tbr{j}aw	&	*pusə\tbr{j}	&	*hua\tbr{j}i	&	*lalə\tbr{j}	&	*qapə\tbr{j}u	&	*pi\tbr{j}a\\ \midrule
Teun	&	{\hr}n-nana	&	{\hr}lou	&	{\hr}lowa	&	{\hr}wulla\su{†}	&	{\hr}wa\tbr{ʔ}i	&	{\hr}lal\tbr{r}a	&	(n)-we\tbr{r}i	&	{\hr}hi\tbr{r}u\\
Nila	&	{\hr}na\tbr{r}na	&	{\hr}lew\tbr{r}a	&	{\hr}le\tbr{r}a	&	{\hr}lus\tbr{r}a	&	{\hr}wa\tbr{r}ni	&	{\hr}lal\tbr{r}a	&	{\hr}mer{\tl}me\tbr{r}u\su{‡}	&	{\hr}i\tbr{r}a\\
Serua	&	{\hr}a\tbr{r}na	&	{\hr}le\tbr{r}a	&	{\hr}le\tbr{r}a	&	{\hr}us\tbr{r}a	&	{\hr}a\tbr{r}ni	&	{\hr}lala\tbr{r}	&	{\hr}mir{\tl}mi\tbr{r}u\su{‡}	&	{\hr}i\tbr{r}a\\
		\lspbottomrule
	\end{tabular}
		\begin{tablenotes}\tnsize
			\item[†]	{Deriving Teun \it{wulla} from *pusəj requires irregular *s
								> \it{l.} The underlying form may be **wulan with /n/ \it{→}
								/l/ /l{\gap} (see note to \trf{04-39}) or perhaps **wulal.
								Assimilation does not affect /lr/, (e.g.
								*laləj > \it{lalra} ‘fly’,
								and \it{falra} ‘wild pigeon’).
								Thus, while this form may have *j > \it{n} or *j > \it{l} it
								does not have *j > \it{r}.}
			\item[‡]	{Nila \it{mer{\tl}meru}
								and Serua \it{mir{\tl}miru} are from *ma-pəju.}
		\end{tablenotes}
	\end{threeparttable}
\end{adjustbox}
\end{table}

Examples of *R are given in \trf{04-42}, \trf{04-43}, and \trf{04-44}.
*maRuqanay has *R > {\0} in both Teun
and Serua and *R > \it{r {\tl}} {\0} in Nila (\trf{04-35}).
One etymon is ambiguous between *R
and *j: *hulaR ‘snake’ or *quləj ‘type of small worm’ > Teun
\it{wula} ‘snake, grub’, Nila \it{laʔulra} ‘snake, grub’, Serua \it{laʔula} ‘snake’.

\begin{table}\tcp{\tsc{Teun-Nila-Serua} *R > \it{r}}{04-42}
\begin{adjustbox}{width=\textwidth}
	\st{2pt}\begin{tabular}{lllllllll}\lsptoprule
local gl.	&	ribs	&	descend	&	conch	&	oyster	&	red	&	new	&	vein	&	wash\\
PMP	&	*\tbr{R}usuk	&	*tu\tbr{R}un	&	*tabu\tbr{R}i	&	*ti\tbr{R}əm	&	*ma-i\tbr{R}aq	&	*baqə\tbr{R}u	&	*u\tbr{R}at	&	*bu\tbr{R}iq\\ \midrule
Teun	&	{\hr}\tbr{r}usi	&	{\hr}tu\tbr{r}un	&	{\hr}sifu\tbr{r}i	&	{\hr}ti\tbr{r}na	&	{\hr}mer{\tl}me\tbr{r}a	&	{\hr}fer{\tl}fe\tbr{r}u	&	{\hr}wu\tbr{r}ʔa	&	{\hr}fu\tbr{r}\\
Nila	&	{\hr}\tbr{r}usʔu	&	{\hr}tu\tbr{r}una	&	{\hr}sifu\tbr{r}i	&	{\hr}ti\tbr{r}mu	&	{\hr}mer{\tl}me\tbr{r}a	&	{\hr}fer{\tl}fe\tbr{r}u	&	{\hr}wu\tbr{r}ta	&	{\hr}fu\tbr{r}\\
Serua	&	\cg{(n-kuʔi)}	&	{\hr}n-tu\tbr{r}un	&	{\hr}miti\tbr{r}i	&	\cg{(tuameti)}	&	(n-)mer{\tl}me\tbr{r}a	&	{\hr}wir{\tl}wi\tbr{r}u	&	{\hr}u\tbr{r}ta	&	{\hr}n-wu\tbr{r}\\
		\lspbottomrule
	\end{tabular}
\end{adjustbox}
\end{table}

\begin{table}\tcp{\tsc{Teun-Nila-Serua} *R /{\gap\#}}{04-43}
\begin{adjustbox}{width=\textwidth}
	\begin{threeparttable}\st{2pt}
	\begin{tabular}{llllll|llll}\lsptoprule
gloss	&	calcium	&	water	&	east	&	testicle	&	root	&	egg	&	coconut	&	tail	&	sail\\
PMP	&	*qapu\tbr{R}	&	*wahiyə\tbr{R}	&	*timu\tbr{R}	&	*lasə\tbr{R}	&	*waka\tbr{R}	&	*qatəlu\tbr{R}	&	*niu\tbr{R}	&	*iku\tbr{R}	&	*laya\tbr{R}\\ \midrule
Teun	&	{\hr}wa\tbr{r}i	&	{\hr}we\tbr{r}u	&	{\hr}tim\tbr{r}i	&	{\hr}las\tbr{r}a	&	{\hr}wow\tbr{r}a\su{†}	&	{\hr}tola	&	{\hr}nowa	&	{\hr}n-iu	&	{\hr}lawa\\
Nila	&	{\hr}wa\tbr{r}u	&	{\hr}we\tbr{r}u	&	{\hr}tim\tbr{r}u	&	{\hr}las\tbr{r}a	&	{\hr}w\<y\>a\tbr{r}a	&	{\hr}tel\tbr{r}u	&	{\hr}nu\tbr{r}u	&	{\hr}n-i\tbr{r}u	&	{\hr}la\tbr{r}u\\
Serua	&	\hp{*w}a\tbr{r}u	&	\hp{*w}e\tbr{r}u	&	{\hr}tim\tbr{r}u	&	{\hr}las\tbr{r}a	&	{\hr}n-e\tbr{r}a	&	{\hr}til\tbr{r}u	&	{\hr}nu\tbr{r}u	&	{\hr}n-i\tbr{r}u	&	{\hr}la\tbr{r}u\\
		\lspbottomrule
	\end{tabular}
		\begin{tablenotes}\tnsize
			\item[†] {\citet[412]{Taber1993} has Teun \it{wowa} ‘root’ which has *R > {\0}.}
		\end{tablenotes}
	\end{threeparttable}
\end{adjustbox}
\end{table}

\begin{table}\tcp{Teun *R > {\0}}{04-44}
\begin{adjustbox}{width=\textwidth}
	\st{2pt}\begin{tabular}{lllllllll}\lsptoprule
gloss	&	house	&	octopus	&	west	&	blood	&	stand	&	charcoal	&	heavy	&	bone\\
PMP	&	*\tbr{R}umaq	&	*ku\tbr{R}ita	&	*haba\tbr{R}at	&	*da\tbr{R}aq	&	*ma-di\tbr{R}i	&	*ba\tbr{R}ah	&	*ba\tbr{R}əqat	&	*du\tbr{R}i\\ \midrule
Teun	&	{\hr}woma	&	{\hr}woʔa	&	{\hr}fata	&	{\hr}rawa	&	{\hr}maruu	&	{\hr}wainfaa	&	{\hr}fo{\tl}fota	&	{\hr}rou\\
Nila	&	{\hr}\tbr{r}uma	&	{\hr}\tbr{r}ita	&	{\hr}fa\tbr{r}ta	&	{\hr}ra\tbr{r}a	&	{\hr}mu-mri\tbr{r}i	&	{\hr}waitfa\tbr{r}a	&	{\hr}fe\tbr{r}ta	&	\cg{(kuʔi)}\\
Serua	&	\cg{(kresna)}	&	{\hr}\tbr{r}ita	&	{\hr}wa\tbr{r}ta	&	{\hr}ra\tbr{r}a	&	{\hr}ne-mri\tbr{r}i	&		&	{\hr}ni-fo\tbr{r}ta	&	\cg{(n-kuʔi)}\\
		\lspbottomrule
	\end{tabular}
\end{adjustbox}
\end{table}

There are two solutions to the reflexes of *d,
*j, and *R in \tsc{Teun-Nila-Serua}.
Firstly, we can posit a merger of *d/*j/*R > \it{r} in Serua and Nila.
This would conflict with the sound changes unifying Teun
and Nila (*ay > \it{a,} *b > \it{f},
{\0} > \it{w} /{\#}{\gap}V),
but align with the lexical data in which Nila
and Serua are closer to one another than either is to Teun.
Alternately, we can posit shared *R/*j > **r at the level of
Proto-TNS, which was followed by **r > {\0},
\it{r} in Teun, after which *d > \it{r} spread by diffusion between the three languages.
These two solutions are outlined in \qf{04-24} and \qf{04-25} below.

{\wg\st{6pt}\tblr{>{\hspace{-\tabcolsep}}l<{\hspace{-\tabcolsep}}>{\hspace{-\tabcolsep}\enspace}lllll}
\tbx{04-24}	&	\mc{4}{l}{Solution 1}\\
						&a.	&*d/*R/*j	&>&\it{r}				&Nila and Serua\\
						&b.	&*R/*j		&>&{\0}, \it{r}	&Teun\\
						&c.	&*d				&>&\it{r}				&Teun\\
\tbx{04-25}	&	\mc{4}{l}{Solution 2}\\
						&a.	&*R/*j		&>&**r					&Teun, Nila, and Serua\\
						&b.	&**r			&>&{\0}, \it{r}	&Teun\\
						&c.	&*d				&>&\it{r}				&Teun, Nila, and Serua\\
\end{tabular}\mbox{}\\}

Whichever solution to the reflexes of *d/*R/*j we adopt,
the lexical similarity and other sound changes indicate there
was a shared period of development and/or contact between Nila
and Teun, as well as Nila
and Serua, but not shared between Serua
and Teun (apart from Nila).
This is unsurprising given the central location of Nila
Island between Teun and Serua Islands (see \frf{04-01}.

